\clearpage
\fancyhead[L]{\small EFFEMERIDI ASTRONOMICHE.}
\SourceLeaf{561}
\begin{center}\small Göttingische gelehrte Anzeigen. Stück 95. Seite 945..952. 14. Juni 1813.\\[3pt]\rule{.72\linewidth}{.3pt}\end{center}

\emph{Effemeridi astronomiche di Milano per l'anno bisestile} 1812 \emph{calcolate da} \textsc{Francesco Carlini e Carlo Brioschi}. \emph{Con appendice.} 1811. Della reale stamperia. Milano. Der Kalender umfasst 124 Seiten, die Zusätze ebenso viele, im Kleinquartformat.

\emph{Effemeridi astronomiche di Milano per l'anno} 1813 usw. 1812. Der Kalender umfasst 95 Seiten, die Zusätze 136 Seiten, im Kleinquartformat.

Die Mailänder astronomischen Ephemeriden sind bei uns wenig bekannt, obwohl sie sowohl durch ihren vorbildlichen Aufbau als auch durch ihre gehaltvollen Zusätze allen vergleichbaren, derzeit erscheinenden Werken Konkurrenz machen. Schade nur, dass sie stets erst kurz vor Beginn des Jahres erscheinen, für das sie bestimmt sind, und deshalb bei uns gewöhnlich erst im Laufe dieses Jahres eintreffen. Wenn sich die Herausgeber doch in dieser Hinsicht Herrn \textsc{Bode} zum Vorbild nähmen! Er lässt seine Jahrbücher immer mehr als zwei Jahre im Voraus erscheinen.

Der Kalender ist folgendermaßen aufgebaut. Neben einem Schlüssel zu den Abkürzungen und einigen Hinweisen zur Zeitrechnung und zu den Kirchenfesten stehen am Anfang eine allgemeine Übersicht der Finsternisse, die scheinbare Schiefe der Ekliptik und die Nutation der Äquinoktialpunkte in Länge, die letzten beiden jeweils einschließlich der Solarnutation. Dann folgen die einzelnen Monate mit je 6 Seiten. Die erste Seite gibt die Mondphasen, die Konjunktionen des Mondes mit Fixsternen, die tatsächlichen Bedeckungen dieser Sterne durch den Mond, enge Konjunktionen der Planeten mit Fixsternen und weitere bemerkenswerte Zeitpunkte im Lauf der Planeten und der Sonne an. Schließlich stehen dort die Finsternisse der Jupitermonde in mittlerer Zeit. Die zweite Seite zeigt den Abstand der Tage vom Jahresanfang, die Wochentage, die mittlere Zeit und die Sternzeit am wahren Mittag, die Sternzeit am mittleren Mittag sowie Sonnenaufgang und Sonnenuntergang. Die dritte Seite enthält die mit größter Sorgfalt auf Zehntelsekunden berechnete Länge der Sonne, ihre Rektaszension im Winkelmaß, ihre Deklination und die Logarithmen ihres Abstands von der Erde mit sechs Dezimalstellen. Die Logarithmen sind für jeden einzelnen Tag angegeben, was ein sehr großer Vorzug gerade dieser Ephemeriden ist. Die vierte und fünfte Seite sind dem Mond gewidmet. Sie geben seine Länge, Breite, Horizontalparallaxe und seinen Durchmesser sowohl für Mittag als auch für Mitternacht an, die Zeiten der Kulmination, des Aufgangs und des Untergangs (nur in Zeitminuten) sowie die Deklination beim Meridiandurchgang in Bogenminuten. Auf der sechsten Seite schließlich sind für jeden Tag des Monats die Konfigurationen der Jupitermonde abgebildet. Dann folgen der Sonnenhalbmesser, die Dauer des Meridiandurchgangs der Sonne und die Länge des Mondknotens von 6 zu 6 Tagen durch das ganze Jahr und anschließend die Bewegungen der einzelnen zehn Planeten. Dass diese in einer einzigen Übersicht für das ganze Jahr zusammengestellt sind und nicht nach einzelnen Monaten aufgeteilt werden, finden wir sehr zweckmäßig. Merkur, Venus und Mars füllen je zwei Seiten, da die Angaben für sie von 6 zu 6 Tagen durch das ganze Jahr fortlaufen. Ceres, Pallas, Juno, Vesta, Jupiter, Saturn und Uranus nehmen nur halb so viel Platz ein, weil die Positionen der ersten vier nur für die Zeit ihrer Sichtbarkeit, die der übrigen nur von 12 zu 12 Tagen angegeben sind. Für jeden Planeten werden auf Minuten genau Länge, Breite, Rektaszension (in Zeitminuten), Deklination, Aufgang, Kulmination und Untergang angegeben. Wenn wir noch einen Wunsch äußern dürften, wäre es der, die Rektaszensionen auf Bogenminuten genau anzugeben (was besonders bei den neuen Planeten wichtig ist) und den Abständen der Planeten von der Erde außerdem eine eigene Spalte zu widmen. Diese Angaben wären besonders für die Berechnung der Aberration sowie für die \SourceLeaf{562}Parallaxe, die scheinbaren Durchmesser und andere Zwecke nützlich. Um dafür Platz zu gewinnen, könnte man je nachdem, ob die Planeten abends oder frühmorgens kulminieren, ihren Aufgang oder ihren Untergang weglassen. Auch könnte die Spalte für die Länge etwas schmaler ausfallen, wenn man die Länge nur in Grad und Minuten statt in Zeichen, Grad und Minuten abdruckte. Wir bemerken noch, dass die Rektaszensionen der Pallas im Jahr 1813 ebenso wie die davon abhängigen Kulminationszeiten wegen eines Rechenfehlers sämtlich falsch sind und dass man für die Rektaszensionen die Komplemente zu 42 Stunden nehmen muss. --- Zuletzt werden noch die Bedeckungen von Fixsternen nach den Rechnungen der Florentiner Astronomen mitgeteilt, so wie sie auch jedes Jahr in der Monatl. Correspondenz abgedruckt werden. Außer diesen regelmäßigen Beiträgen sind dem Kalender für 1812 ein Verzeichnis aller in Mailand sichtbaren Sterne über der fünften Größe beigefügt, nach \textsc{Piazzi} auf 181c reduziert (insgesamt 596), und die \textsc{Gauß}schen Tafeln für Aberration und Nutation zusammen mit einer kleinen Tafel für die Solarnutation.

Unter den \emph{Zusätzen} zum Jahrgang 1812 stehen die Beobachtungen der Zenitdistanzen der Sonne und der Fixsterne im Meridian an erster Stelle, die \textsc{Barnabas Oriani} mit einem neuen Wiederholungskreis angestellt hat. Vorangestellt ist eine lehrreiche Beschreibung des unvergleichlichen Instruments, eines drei Fuß großen \textsc{Reichenbach}schen Kreises mit aufrechter Säule (Preis 3000 Gulden). Die mit diesem Instrument über zwei Jahre hinweg angestellten und in den beiden Jahrgängen 1812 und 1813 abgedruckten Beobachtungen enthalten einen Schatz an Erfahrungen für die Sterndeklinationen, für die Theorie der Sonnenbewegung und für die Theorie der astronomischen Strahlenbrechung. Wir wünschen, dass dieser Schatz bald so umfassend genutzt wird, wie er es verdient. --- Beobachtungen zur Bestimmung der Sonnenwenden und der Schiefe der Ekliptik in den Jahren 1810 und 1811, von \textsc{Angelo Cesaris}. Während \textsc{Oriani} die Wintersonnenwende von 1810 und die Sommersonnenwende von 1811 mit dem drei Fuß großen \textsc{Reichenbach}schen Kreis beobachtete, benutzte \textsc{Cesaris} zum gleichen Zweck den achtfüßigen Mauerquadranten. Die hier mitgeteilten Ergebnisse beider Beobachter stimmen auf einige Zehntelsekunden überein. Wir hätten uns nur noch eine Erklärung gewünscht, auf welchem Weg der Kollimationsfehler des Quadranten ermittelt wurde, den \textsc{Cesaris} für beide Sonnenwenden mit $= -1''5$ ansetzt. Auch deuten diese Beobachtungen auf keinen Unterschied der Schiefe zwischen Winter- und Sommersonnenwende hin, wie ihn einige Astronomen gefunden haben wollen. Denn die von \textsc{Cesaris} unter der Annahme gleicher Schiefe abgeleitete Polhöhe $45^\circ 28' 0''20$ stimmt sehr genau mit der auf anderem Weg gefundenen überein. --- Über den Konvergenzgrad der verschiedenen Reihen, die die Ungleichheiten der Mondlänge darstellen, von \textsc{Franz Carlini}. Ein außerordentlich wertvoller Aufsatz! Hier werden drei Ausdrücke für die Mondungleichheiten zusammengestellt. Der erste folgt der Form von \textsc{Tobias Mayer}, nach der die Tafeln von \textsc{Bürg} berechnet sind. Der zweite, von \textsc{Laplace} aus dem ersten abgeleitet, hat die Form, in der die Theorie von \textsc{Laplace} diese Ungleichheiten ergeben hat, also als Funktionen der wahren Mondlänge. Der dritte, von \textsc{Carlini} aus dem ersten berechnet, hat die Form, in die schon \textsc{Lambert} und \textsc{Schulze} die Gleichungen von \textsc{Mayer} und \textsc{Mason} gebracht hatten, also als Funktionen allein der mittleren Bewegungen. Wir würden diese letzte Form allen anderen vorziehen, wenn der Ausdruck auch etwas langsamer konvergiert als die ersten beiden. Es wäre wünschenswert, \emph{alle} seine Koeffizienten unmittelbar aus einigen Tausend Beobachtungen abzuleiten. Das könnten freilich nur die Astronomen mit Vorteil ausführen, denen dabei die Vorarbeiten von \textsc{Bürg} oder \textsc{Burckhardt} zur Verfügung stünden. \textsc{Carlini} hat die Prüfung selbst an 100 aus Sternbedeckungen abgeleiteten Mondpositionen vorgenommen und durchweg gute Übereinstimmung gefunden. Bemerkenswert ist, dass \textsc{Carlini} hier eine neue Gleichung bestätigt, die \textsc{Burckhardt} schon früher aufgestellt, jetzt aber in seine neuen Tafeln nicht aufgenommen hat. Sie kündigt sich zunächst in Form einer Ungleichheit der Exzen\SourceLeaf{563}trizität mit langer Periode an. --- Über den Einfluss von Temperaturänderungen auf die Bewegungen des Pendels, von \textsc{Carl Brioschi}. Der Verfasser hat diese Untersuchung mit größerem methodischem Aufwand durchgeführt, als nötig gewesen wäre. Im Wesentlichen besteht das Ergebnis doch nur darin, dass ein Pendel, dessen Länge sich bei veränderlicher Temperatur ändert, in einem bestimmten Zeitintervall genau ebenso viele Schwingungen ausführt, wie es mit konstanter Länge beim mittleren Thermometerstand ausgeführt hätte. --- Zuletzt noch einige Sternbedeckungen, die \textsc{Carlini} und \textsc{Oriani} in den Jahren 1808 und 1810 in Mailand und Rom beobachtet haben.

Die Zusätze zum Jahrgang 1813 beginnen mit der bereits erwähnten Fortsetzung der mit dem drei Fuß großen \textsc{Reichenbach}schen Kreis beobachteten Zenitdistanzen. --- Über das periodische Schwanken der Gebäude, von \textsc{Angelo Cesaris}. Seit Langem weiß man, wie wenig man sich darauf verlassen kann, dass astronomische Instrumente, die eine feste Aufstellung erfordern, ihre Lage unverändert beibehalten, wenn sie mit den Außenmauern des Gebäudes verbunden sind, besonders in beträchtlicher Höhe über dem Erdboden. Auf diese Mauern wirken Temperatur und Feuchtigkeit unmittelbar ein. Daher ist es wichtig, solchen Instrumenten beim Bau einer neuen Sternwarte zu ebener Erde eine vollkommen feste und von den Außenmauern ganz unabhängige Grundlage zu geben. Die Sternwarte Brera in Mailand, die sonst mit so ausgezeichneten Instrumenten ausgestattet ist, hat diesen Vorzug nicht. Die Schwankungen des Gebäudes zeigen sich am Mauerquadranten und am Passageninstrument sehr deutlich und stark. Wenn der senkrechte Faden des Passageninstruments morgens vor Sonnenaufgang das in 3000 Toisen Entfernung errichtete Meridianzeichen berührt, entfernt sich der Faden anschließend immer weiter nach Osten von diesem Zeichen bis nach Mittag. Dann beträgt der Abstand an heiteren Wintertagen 5 bis 6 Sekunden, an heiteren Sonnentagen aber 30 Sekunden. Unter sonst gleichen Umständen ist die Veränderung bei bedecktem Himmel und nahezu gleichbleibender Temperatur am geringsten oder ganz unmerklich, bei Sonnenschein und starker Temperaturänderung dagegen am größten. --- Über die Formeln für die Parallaxe und die Breite des Mondes, von \textsc{Franz Carlini}, eine Fortsetzung des Aufsatzes im vorhergehenden Band. Ähnlich wie die Formel für die Ungleichheiten der Länge wird die Formel für die Parallaxe in eine andere umgewandelt, die nur von mittleren Bewegungen abhängt. Einen ähnlichen Ausdruck hat \textsc{Carlini} für den Logarithmus der Parallaxe entwickelt. Den Ausdruck für die Breite hat er hingegen nicht auf dieselbe Weise umgeformt, weil er dann viel zu schwerfällig geworden wäre und zu langsam konvergiert hätte. \textsc{Carlini} beschränkt sich darauf, ein sehr einfaches Verfahren anzugeben, mit dem sich Tafeln für die Form von \textsc{Laplace} (in Funktionen der wahren Länge) bequem einrichten lassen. Aus mehreren Gründen hätten wir dennoch gewünscht, dass Carlini auch bei diesem Element, \emph{nur unter Ausschluss der ersten beiden Glieder}
\[
18540''25\sin\delta + 12''56\sin 3\delta\text{,}
\]
diese Umformung vorgenommen hätte. Dann wäre die Schwierigkeit entfallen, und die Analogie zwischen den Formeln für die Ungleichheiten des Mondes und denen für die Ungleichheiten der Planeten wäre dadurch vollständig geworden. Die Umformung lässt sich indessen leicht nachholen. --- Opposition des Saturn im Jahr 1811, beobachtet von \textsc{Carl Brioschi}. --- Auszug aus den meteorologischen Beobachtungen an der Mailänder Sternwarte im Jahr 1809, von \textsc{Angelo Cesaris}.

